% ==== AUTO-GENERADA — no editar a mano (regenerar con generate_thesis_tables.py) ====
\begin{table}[htbp]
\centering
\caption[Perspectiva global de la campaña por topología]{Perspectiva global de la campaña por topología. La columna \emph{Calif.\ mejor punto} califica el punto de menor $|\Delta E|$ alcanzado por la topología en toda la campaña; \emph{$|\Delta E|$ medio} promedia sobre todos los puntos evaluados de todos los $N$ y profundidades de esa topología, incluidos los situados por debajo de la frontera $h_{\min}(N, p)$, y \emph{Calif.\ media} lo califica. Por eso esta tabla mide cobertura de la campaña y su $|\Delta E|$ medio es superior al de la Tabla~\ref{tab:cross_topo_de}, que se restringe a $N = 10$, a la mejor profundidad y al régimen válido $h \geq 1{,}3$ (rendimiento en el régimen de operación). La columna \emph{$N$ máx.} es el mayor tamaño con resultado medido para esa topología en la campaña (no el límite del método, que sigue las tres escalas de la Tabla~\ref{tab:comparison_literature}). Escala de calificación A--E por $|\Delta E|$: A ($<0{,}05$), B ($<0{,}10$), C ($<0{,}30$), D ($<1{,}00$), E ($\geq 1{,}00$); se usa E como peor nota para no confundirla con el símbolo de fidelidad $F$. Fuente: elaboración propia.}
\label{tab:auto_scoreboard}
\begin{tabular}{lcccc}
\toprule
Topología & $N$ máx. & Calif. mejor punto & $|\Delta E|$ medio & Calif. media \\
\midrule
Cadena 1D & 100 & B & 0,1882 & C \\
Heavy-hex & 60 & A & 0,1242 & C \\
Escalera & 40 & B & 0,2451 & C \\
Cuadrada & 30 & B & 0,1588 & C \\
Triangular & 24 & B & 0,1164 & C \\
\bottomrule
\end{tabular}
\end{table}
